
Large language models achieve high code generation accuracy but exhibit lower success rates when attempting to repair their own compiler errors during self-repair loops. This work investigates whether restructuring compiler error messages improves repair outcomes. We introduce structured error formatting---organizing compiler output with explicit section labels (PROBLEM/LOCATION/CONTEXT/ROOT\_CAUSE)---and test whether this organizational structure improves self-repair success rates. Experiments on EvalPlus benchmarks compare structured formatting against a scrambled control condition containing identical information with randomized section order. Results show that structured formatting achieves 48.4\% repair success compared to 34.0\% for scrambled formatting ($\Delta$ = +14.4 percentage points, p < 0.001, McNemar's test), providing evidence that organizational structure affects repair success independently of information content. An information preservation test confirms 100\% reconstruction accuracy across 500 samples, validating that format transformations preserve diagnostic content. A fix specificity experiment conducted in simulation mode shows an inverted-U pattern with peak success at Level 2 (specific patterns, 60.2\%) compared to no hints (34.7\%) and exact fixes (39.6\%). A scale interaction test finds a directional pattern (smaller models show larger format benefits) but the interaction does not reach statistical significance (p = 0.176, $\eta^2$ = 0.001). These findings suggest that error format organization may be a relevant variable in self-repair system design.
